\section{Esempio di Istanze di Test}
\label{sec:app_abox}

Per la verifica end-to-end del sistema di reasoning, \`e stato utilizzato un
insieme di istanze di test minimale (\texttt{test\_abox.ttl}) contenente 3 giochi, 2 developer
e alcune piattaforme. Il seguente estratto in Turtle illustra la struttura:

\begin{lstlisting}[style=turtle, caption={Estratto del test delle istanze (Turtle)}, label={lst:testabox}]
@prefix vg: <http://www.videogame-ontology.org/ontology#> .

vg:GameA  a vg:VideoGame ;
    vg:gameName "Elden Ring" ;
    vg:releaseDate "2022-02-25"^^xsd:date ;
    vg:metacriticScore 96 ;
    vg:developedBy vg:Dev1 ;
    vg:hasGenre vg:RPGGenreInstance ;
    vg:hasGameMode vg:SinglePlayer, vg:OnlineMultiplayer ;
    vg:availableOn vg:PS5, vg:PC_Windows ;
    vg:wonAward vg:TGAAward2022 .

vg:GameB  a vg:VideoGame ;
    vg:gameName "Hollow Knight" ;
    vg:releaseDate "2017-02-24"^^xsd:date ;
    vg:metacriticScore 87 ;
    vg:developedBy vg:Dev2 ;
    vg:availableOn vg:NintendoSwitch ;
    vg:hasGameMode vg:SinglePlayer .

vg:GameC  a vg:VideoGame ;
    vg:gameName "Dark Souls III" ;
    vg:sequelOf vg:GameX ;
    vg:developedBy vg:Dev1 .

# individui tipizzati (necessari per le restrizioni esistenziali)
vg:PS5               a vg:PlayStationPlatform .
vg:PC_Windows        a vg:WindowsPlatform .
vg:NintendoSwitch    a vg:NintendoPlatform .
vg:SinglePlayer      a vg:SinglePlayerMode .
vg:OnlineMultiplayer a vg:MultiplayerMode .
vg:TGAAward2022      a vg:IndustryAward .
\end{lstlisting}

\subsection{Risultati della Verifica}

Dopo il reasoning completo (owlrl + CONSTRUCT), il grafo di test produce:

\begin{itemize}[nosep]
  \item \textbf{2.715 triple finali} (dalle $\sim$997 iniziali);
  \item \textbf{0 \texttt{owl:sameAs} falsi} --- il fix degli assiomi DL-only
    previene l'unificazione errata dei giochi;
  \item \textbf{Classificazioni corrette}:
    \begin{itemize}[nosep]
      \item \texttt{GameA} $\rightarrow$ MultiplayerGame, SinglePlayerGame,
        ConsoleGame, PlayStationGame, PCGame, HighlyRatedGame, RecentGame,
        AwardWinningGame, StandaloneGame;
      \item \texttt{GameB} $\rightarrow$ SinglePlayerGame, ConsoleGame,
        NintendoGame, HighlyRatedGame, StandaloneGame (non RecentGame: 2017);
      \item \texttt{GameC} $\rightarrow$ SequelGame (nessuna classificazione per
        piattaforma o modalit\`a: non sono asseriti \texttt{availableOn} n\'e
        \texttt{hasGameMode}).
    \end{itemize}
  \item \textbf{Propriet\`a derivate}:
    \begin{itemize}[nosep]
      \item \texttt{Dev1} \texttt{developerOf} \texttt{GameA}, \texttt{GameC}
        (per l'inversa di \texttt{developedBy});
      \item \texttt{GameA} \texttt{sharesDeveloperWith} \texttt{GameC} e viceversa
        (catena \texttt{developedBy} $\circ$ \texttt{developerOf} e simmetria);
      \item \texttt{GameX} \texttt{hasSequel} \texttt{GameC} (inversa); \texttt{GameX}
        \`e inoltre classificato \texttt{VideoGame} per effetto del range di
        \texttt{sequelOf}.
    \end{itemize}
\end{itemize}
